\section{Вывод}

В лабораторной работе исследовано оценивание математического ожидания гауссовского случайного процесса и влияние параметров выборки на точность оценки. По реализациям и спектральным плотностям мощности установлено, что с ростом времени корреляции характерные изменения сигнала становятся более медленными, а спектр сужается.

Эксперименты с зависимостью СКО от числа отсчётов показали уменьшение ошибки по мере увеличения $N$. Увеличение интервала $\Delta t$ снижает корреляцию между используемыми отсчётами и повышает эффективность усреднения. Для достаточно больших интервалов СКО приближается к пределу, соответствующему усреднению практически независимых отсчётов.

В задании~6 установлено, что доверительный интервал уменьшается при росте числа усредняемых отсчётов и увеличивается при повышении доверительной вероятности. Полученные таблицы и гистограммы подтверждают эти закономерности.

Спектральные оценки, полученные в задании~5, согласуются с результатами обработки временных реализаций: среднее по одной реализации отличается от оценки по ансамблю в пределах статистического разброса, а площадь под непрерывной частью спектра соответствует дисперсии процесса.

Таким образом, точность оценки среднего зависит от числа усредняемых отсчётов, интервала между ними и корреляционных свойств процесса.
